% Software Design Description — Vehicle Dynamics (IEEEtran)
% Requires IEEEtran.cls (TeX Live / MiKTeX IEEE package).

\documentclass[conference]{ieeetran.cls}
\usepackage{amsmath,amssymb,amsfonts}
\usepackage{booktabs}
\usepackage{cite}
\usepackage{url}
\usepackage{hyperref}

\begin{document}

\title{Software Design Description\\for a Vehicle Dynamics Simulator with DEM Soft Soil}

\author{\IEEEauthorblockN{Sankalp Amaravadi}
\IEEEauthorblockA{Vehicle Dynamics Project\\
Email: (TBD)}}

\maketitle

\begin{abstract}
This Software Design Description (SDD) outlines the architecture of a vehicle
dynamics backend that couples chassis/tire models to GPU-oriented discrete
element method (DEM) soft-soil simulation. The design separates orchestration,
vehicle physics, granular soil, tire--soil coupling, and I/O. Implementation
details for individual modules remain TBD pending Phase~2 development.
\end{abstract}

\begin{IEEEkeywords}
software design, vehicle dynamics, DEM, architecture, GPU, Python
\end{IEEEkeywords}

\section{Introduction}

\subsection{Purpose}
This SDD communicates system structure, module boundaries, and planned data
flow for implementers and reviewers.

\subsection{Scope}
Covered: Python orchestration, vehicle subsystem, DEM soil subsystem, coupling,
integration loop, and export. Not covered: graphical front end.

\subsection{Reference Documents}
Companion artifacts include the project SRS (\texttt{srs.tex}) and testing
document (\texttt{testing.tex}) under \texttt{documentation/}.

\section{System Overview}
The simulator is a greenfield application with entry point \texttt{main.py}.
Backend packages are planned under \texttt{src/}. Documentation is maintained as
IEEE-formatted \LaTeX\ sources.

\section{Architectural Design}

\subsection{Logical View}
\begin{enumerate}
  \item \textbf{Orchestration} --- configuration loading, scenario selection,
        and the global time loop.
  \item \textbf{Vehicle} --- multi-DOF chassis and tire models (beyond
        bicycle-only approximations).
  \item \textbf{DEM soil} --- particle state, neighbor search, and contact
        resolution; GPU kernels for large $N$.
  \item \textbf{Coupling} --- exchange of forces and kinematic constraints
        between wheels and soil particles.
  \item \textbf{I/O} --- logging and CSV/time-series export.
\end{enumerate}

\subsection{Process View}
A fixed- or adaptive-step loop (TBD) advances vehicle and DEM state, then
applies coupling forces before optional export at configured intervals.

\subsection{Data Flow}
Configuration $\rightarrow$ initialize vehicle and particle field $\rightarrow$
time step (vehicle, DEM, couple) $\rightarrow$ export.

\section{Detailed Design (TBD)}
Class diagrams, interfaces, and algorithm notes for DEM contact and vehicle
integrators will be added as modules are implemented.

\section{Deployment View}
Primary deployment is a local workstation. Million-scale DEM runs expect an
NVIDIA-compatible GPU and drivers. Operators invoke the system via
\texttt{python main.py}.

\section*{Acknowledgment}
This document follows the IEEE conference \LaTeX\ template (\texttt{IEEEtran}).

\begin{thebibliography}{00}
\bibitem{ieee1016} IEEE Std 1016-2009, \emph{IEEE Standard for Information
Technology---Systems Design---Software Design Descriptions}.
\bibitem{ieeetran} M.~Shell, ``IEEEtran homepage,''
\url{https://www.michaelshell.org/tex/ieeetran/}.
\end{thebibliography}

\end{document}
